\documentclass[11pt]{article}

\usepackage[margin=1in]{geometry}
\usepackage{enumitem}
\usepackage{titlesec}
\usepackage{hyperref}
\usepackage{xcolor}
\usepackage{tabularx}
\usepackage{booktabs}
\usepackage{longtable}

\titleformat{\section}{\large\bfseries}{}{0em}{}[\titlerule]
\titleformat{\subsection}{\normalsize\bfseries}{}{0em}{}

\hypersetup{colorlinks=true, linkcolor=blue, urlcolor=blue}

\begin{document}

% ============================================================
% HEADER
% ============================================================
\begin{center}
    {\LARGE \textbf{Sprint Plan: Sprint \#3}} \\[6pt]
    {\large CS161, Spring 2026}
\end{center}

\vspace{0.5em}

% ============================================================
% 1. Team Information
% ============================================================
\section{Team Information}

\begin{tabularx}{\textwidth}{@{}lX@{}}
    \textbf{Team Number:}          & Team 2 \\
    \textbf{Team Name:}            & Checkpoint \\
    \textbf{Team Members:}         & Omkar Podey, Anmol Rastogi \\
    \textbf{Class:}                & CS161, Software Engineering, Spring 2026 \\
    \textbf{Planning Date:}        & April 9, 2026 \\
    \textbf{Sprint Period:}        & April 10 to April 23, 2026 (2 weeks) \\
\end{tabularx}

% ============================================================
% 2. Task List Link
% ============================================================
\section{Product Backlog / Task List}

\textbf{Trello Board:} \url{https://trello.com/b/ufZ3LwJi/checkpoint-cs161}

\textbf{GitHub Repository:} \url{https://github.com/Omkar399/checkpoint}

% ============================================================
% 3. Sprint 3 Context
% ============================================================
\section{Sprint Context: Building on Sprint 2}

Sprint~1 established the foundational chat infrastructure (servers, channels, real-time messaging). Sprint~2 delivered the core accountability feature (check-ins, daily dashboard, heatmaps). Sprint~3 adds the social and automated engagement layer: the Coach Bot for smart notifications, leaderboards for peer recognition, emoji reactions for quick feedback, and testing infrastructure for maintainability.

% ============================================================
% 4. Sprint 3 Objectives
% ============================================================
\section{Sprint 3: Objectives}

The primary objective of Sprint~3 is to enhance user engagement through automated coaching, social recognition, and continuous quality improvement.

\textbf{Key goals:}
\begin{itemize}[leftmargin=1.5em]
    \item Implement the Coach Bot: automated welcome messages, daily summaries, and inactivity nudges.
    \item Add emoji reaction support to check-in cards for peer encouragement.
    \item Build a Leaderboard component showing monthly rankings by check-in count.
    \item Polish the UI with loading states, error toasts, empty states, and responsive design.
    \item Begin writing automated API tests (at least 20 tests covering auth, check-ins, and channels).
    \item Improve code maintainability with test fixtures and shared utilities.
\end{itemize}

% ============================================================
% 5. Sprint 3 User Stories
% ============================================================
\section{Sprint 3: User Stories}

\begin{center}
\begin{tabularx}{\textwidth}{@{}clXc@{}}
\toprule
\textbf{ID} & \textbf{Priority} & \textbf{Story} & \textbf{Status} \\
\midrule
3.1 & Must Have & As a new channel member, I receive a Coach Bot welcome message explaining the check-in system. & \textbf{Planned} \\
3.2 & Must Have & As a channel member, I receive a daily summary of that day's check-ins at 9 PM. & \textbf{Planned} \\
3.3 & Should Have & As a channel member, I receive a 48-hour inactivity nudge if I haven't checked in in two days. & \textbf{Planned} \\
3.4 & Must Have & As a channel member, I can react to check-in cards with emoji (thumbs up, heart, fire, etc.). & \textbf{Planned} \\
3.5 & Must Have & As a channel member, I can view a monthly leaderboard showing peers ranked by check-in count. & \textbf{Planned} \\
3.6 & Should Have & As a UI user, I see loading skeletons, error toasts, and empty state illustrations. & \textbf{Planned} \\
3.7 & Must Have & As a developer, there are at least 20 automated tests for the API. & \textbf{Planned} \\
\bottomrule
\end{tabularx}
\end{center}

\textbf{Rationale:} These stories address user retention (Coach Bot, nudges), peer encouragement (reactions, leaderboard), and codebase health (tests, UI polish).

% ============================================================
% 6. Sprint 3 Task Breakdown
% ============================================================
\section{Sprint 3: Planned Task Breakdown}

All 7 Sprint~3 tasks are planned to be completed by April 23, 2026.

\begin{center}
\begin{tabularx}{\textwidth}{@{}lXcc@{}}
\toprule
\textbf{\#} & \textbf{Task} & \textbf{Assigned To} & \textbf{Est. Days} \\
\midrule
1 & Coach Bot: Welcome, Daily Summary, Inactivity Nudges & Omkar & 2 \\
2 & Emoji Reactions: Model, Service, API, Storage & Omkar & 1.5 \\
3 & Leaderboard Component \& Integration & Anmol & 1.5 \\
4 & UI Polish: Loading States, Toasts, Empty States & Anmol & 1.5 \\
5 & Emoji Reactions Frontend: Display, Interaction & Anmol & 1 \\
6 & API Testing Suite: 20+ Tests (auth, check-ins, channels, leaderboard) & Omkar & 2 \\
7 & Integration \& Final QA & Both & 1 \\
\bottomrule
\end{tabularx}
\end{center}

% ============================================================
% 7. Technical Architecture: Sprint 3
% ============================================================
\section{Technical Architecture: Sprint 3}

\subsection{Coach Bot (Backend)}

The Coach Bot is a scheduled background task that sends system messages to channels:

\begin{itemize}[leftmargin=1.5em]
    \item \textbf{Welcome Message}: Sent when a user joins a channel, explaining the check-in ritual and the purpose of the daily dashboard.
    \item \textbf{Daily Summary}: Sent at 9 PM (UTC) listing all members and their check-in status for that day.
    \item \textbf{Inactivity Nudge}: Sent at 9 AM if the user hasn't checked in for 48 hours, encouraging them to return.
\end{itemize}

\textbf{Implementation:}
\begin{itemize}[leftmargin=1.5em]
    \item New \texttt{CoachMessage} SQLAlchemy model (id, channel\_id, user\_id, message\_type, sent\_at).
    \item New \texttt{coach\_service} with functions: \texttt{send\_welcome}, \texttt{send\_daily\_summary}, \texttt{send\_inactivity\_nudge}.
    \item APScheduler (or similar) background worker for scheduled tasks.
    \item Broadcast coach messages to all connected channel members via WebSocket.
\end{itemize}

\subsection{Emoji Reactions (Backend)}

Add support for emoji reactions on check-in cards:

\begin{itemize}[leftmargin=1.5em]
    \item New \texttt{Reaction} model (id, checkin\_id, user\_id, emoji, created\_at).
    \item New endpoints:
    \begin{itemize}
        \item POST \texttt{/checkins/\{id\}/reactions} --- Add a reaction.
        \item DELETE \texttt{/checkins/\{id\}/reactions/\{emoji\}} --- Remove a reaction.
        \item GET \texttt{/checkins/\{id\}/reactions} --- List all reactions (aggregated by emoji with counts).
    \end{itemize}
    \item WebSocket event: Broadcast \texttt{new\_reaction} when a reaction is added/removed.
\end{itemize}

\subsection{Leaderboard (Backend \& Frontend)}

Monthly ranking of members by check-in count:

\begin{itemize}[leftmargin=1.5em]
    \item New endpoint: GET \texttt{/channels/\{id\}/leaderboard?year=2026\&month=4} --- Returns list of (user, check-in\_count, rank).
    \item Service function: \texttt{get\_leaderboard(db, channel\_id, month, year)} --- Aggregates check-in counts for the given month.
    \item Frontend component: \texttt{Leaderboard.jsx} --- Table/list showing rank, user avatar, username, check-in count, and badges for top 3.
\end{itemize}

\subsection{Testing Infrastructure}

\begin{itemize}[leftmargin=1.5em]
    \item Use \texttt{pytest} with fixtures for:
    \begin{itemize}
        \item Database session (in-memory SQLite for tests).
        \item Authenticated client (test user with valid JWT).
        \item Test server, channel, and membership setup.
    \end{itemize}
    \item Target test coverage:
    \begin{itemize}
        \item Authentication: login, register, token refresh, unauthorized access.
        \item Check-ins: create, list, dashboard, streak, heatmap.
        \item Channels: create, invite, member list, not-a-member 403.
        \item Leaderboard: fetch, sorting, monthly filters.
    \end{itemize}
    \item All tests should be runnable via \texttt{pytest} with minimal setup.
\end{itemize}

\subsection{UI Polish}

\begin{itemize}[leftmargin=1.5em]
    \item Loading states: Skeleton loaders for MessageFeed, DailyDashboard, Leaderboard.
    \item Error handling: Toast notifications (top-right, auto-dismiss) for API failures.
    \item Empty states: Illustrations and helpful text when no messages, no members, or no check-ins exist.
    \item Responsive design: Test layouts on mobile (375px), tablet (768px), and desktop (1920px).
\end{itemize}

% ============================================================
% 8. Implementation Strategy
% ============================================================
\section{Implementation Strategy}

\subsection{Backend-First Approach}

1. \textbf{Coach Bot} (Day 1-2): Implement model, service, and background scheduler. Test message generation logic.
2. \textbf{Emoji Reactions} (Day 1.5-2.5): Implement model, endpoints, and WebSocket broadcast. Keep simple: 7 predefined emoji.
3. \textbf{Leaderboard Service} (Day 2-3): Implement database query and REST endpoint. This unblocks frontend.
4. \textbf{API Tests} (Day 4-5): Write comprehensive test suite once backend features are stable.

\subsection{Frontend-Parallel Approach}

1. \textbf{Emoji Reactions UI} (Day 2-3): Once backend endpoints are available, build reaction button and reaction list.
2. \textbf{Leaderboard Component} (Day 2-3): Fetch leaderboard data, render table with sorting and badges.
3. \textbf{UI Polish} (Day 3-4): Add loading skeletons, toast components, empty states, and responsive adjustments.

\subsection{Integration \& QA (Day 5)}

1. End-to-end testing: Post check-in $\rightarrow$ Receive reaction $\rightarrow$ View leaderboard.
2. Coach Bot message delivery verification.
3. Mobile responsiveness check.

% ============================================================
% 9. Risk Mitigation
% ============================================================
\section{Risk Mitigation}

\begin{enumerate}[leftmargin=1.5em]
    \item \textbf{Risk:} Background scheduler (APScheduler) is a new external dependency.
    \begin{itemize}
        \item \textbf{Mitigation:} Use APScheduler v3.x (stable). Implement a simple fallback: manual background task triggering via a GET endpoint for testing.
    \end{itemize}

    \item \textbf{Risk:} Reactions and leaderboard queries could be slow with large data.
    \begin{itemize}
        \item \textbf{Mitigation:} Add database indexes on (checkin\_id) for reactions and (channel\_id, checked\_in\_at) for leaderboard aggregation (already planned from Sprint~2).
    \end{itemize}

    \item \textbf{Risk:} Writing 20+ tests is time-consuming with limited sprint time.
    \begin{itemize}
        \item \textbf{Mitigation:} Prioritize core auth and check-in tests first (aim for 12 tests by Day 4). Additional tests can be added in a post-sprint refactoring session.
    \end{itemize}

    \item \textbf{Risk:} UI polish can scope-creep if not time-boxed.
    \begin{itemize}
        \item \textbf{Mitigation:} Define specific polish items: 3 loading skeleton components, 1 toast component, 3 empty state illustrations. Limit time to 1.5 days.
    \end{itemize}
\end{enumerate}

% ============================================================
% 10. Definition of Done (Sprint 3)
% ============================================================
\section{Definition of Done: Sprint 3 Criteria}

\begin{enumerate}[leftmargin=1.5em]
    \item \textbf{All 7 user stories (3.1--3.7) are implemented and manually tested.}

    \item \textbf{Coach Bot successfully sends welcome, daily summary, and inactivity nudge messages.}

    \item \textbf{Emoji reactions appear in real-time on check-in cards and persist in the database.}

    \item \textbf{Leaderboard component displays correctly and handles monthly filtering.}

    \item \textbf{At least 15 automated tests pass (covering auth, check-ins, channels, and leaderboard).}

    \item \textbf{UI includes loading states, error toasts, and empty state illustrations.}

    \item \textbf{Mobile responsive design verified on 375px, 768px, and 1920px viewports.}

    \item \textbf{All code follows existing patterns and passes a final code review.}
\end{enumerate}

% ============================================================
% 11. Dependencies and Blockers
% ============================================================
\section{Dependencies and Blockers}

\begin{enumerate}[leftmargin=1.5em]
    \item \textbf{APScheduler Installation:} Required for background task scheduling. No blocker, as it's a standard package.

    \item \textbf{React Toast Component:} Can use existing community libraries (e.g., \texttt{react-hot-toast}) or build a simple custom component. Decision needed by Day 3.

    \item \textbf{Emoji Set:} Define the 7 allowed emoji early (suggested: \texttt{+1}, \texttt{heart}, \texttt{fire}, \texttt{laughing}, \texttt{thinking}, \texttt{tada}, \texttt{rocket}).
\end{enumerate}

% ============================================================
% 12. Success Metrics
% ============================================================
\section{Success Metrics}

\begin{enumerate}[leftmargin=1.5em]
    \item \textbf{Velocity:} Complete all 7 user stories and tasks on time (by April 23).

    \item \textbf{Test Coverage:} Achieve 15+ passing tests and 60\%+ code coverage for backend API.

    \item \textbf{User Engagement:} Platform now has three engagement layers: check-ins (core), reactions (quick feedback), leaderboard (recognition).

    \item \textbf{Code Quality:} No critical bugs reported in end-to-end testing; UI is responsive and polished.

    \item \textbf{Technical Debt Reduction:} Establish testing practices for future sprints.
\end{enumerate}

% ============================================================
% 13. Sprint Schedule
% ============================================================
\section{Sprint Schedule}

\begin{center}
\begin{tabularx}{\textwidth}{@{}clXc@{}}
\toprule
\textbf{Day} & \textbf{Date} & \textbf{Focus} & \textbf{Deliverable} \\
\midrule
1 & Apr 10 (Thu) & Coach Bot model \& service & Welcome \& daily summary logic \\
2 & Apr 11 (Fri) & Emoji Reactions model \& endpoints & Reactions CRUD endpoints \\
3 & Apr 14 (Mon) & Leaderboard service \& frontend & Leaderboard component working \\
4 & Apr 15 (Tue) & Emoji reactions frontend & Reaction button \& display in UI \\
5 & Apr 16 (Wed) & UI polish \& loading states & Skeletons, toasts, empty states \\
6 & Apr 17 (Thu) & API tests (batch 1) & 15+ passing tests \\
7 & Apr 18 (Fri) & Integration \& final testing & End-to-end verification \\
8 & Apr 21 (Mon) & Buffer \& refinement & Hotfixes, final review \\
9 & Apr 22 (Tue) & Final QA \& documentation & Report generation \\
10 & Apr 23 (Wed) & Sprint review \& handoff & Final deliverables \\
\bottomrule
\end{tabularx}
\end{center}

Scrum meetings: 2 per week (in-class + virtual).

% ============================================================
% 14. Rollover / Scope Management
% ============================================================
\section{Scope Management}

\textbf{If time is running short, the priority order is:}

\begin{enumerate}[leftmargin=1.5em]
    \item \textbf{P1 (Must Complete):} Coach Bot (3.1, 3.2), Emoji Reactions (3.4), Leaderboard (3.5).
    \item \textbf{P2 (Should Complete):} Inactivity Nudges (3.3), API Tests (3.7).
    \item \textbf{P3 (Nice-to-Have):} UI Polish details (3.6), additional tests beyond 15.
\end{enumerate}

If tasks 3.3 or 3.7 roll over, they will be prioritized in Sprint~4.

% ============================================================
% 15. Next Steps (Post-Sprint 3)
% ============================================================
\section{Sprint 4: Preliminary Ideas}

Based on the roadmap in the PRD, Sprint~4 could include:

\begin{enumerate}[leftmargin=1.5em]
    \item \textbf{Goals / Habit Tracking:} Long-term challenge tracking (e.g., ``Run 100 miles in a month'').
    \item \textbf{Export / Reporting:} CSV export of monthly check-in summaries.
    \item \textbf{Mobile App:} React Native companion app for iOS/Android.
    \item \textbf{Analytics Dashboard:} Server-level insights (total check-ins, member engagement over time).
    \item \textbf{Custom Channels:} Template-based channels with predefined check-in types (e.g., fitness, productivity, learning).
\end{enumerate}

% ============================================================
% 16. Glossary \& References
% ============================================================
\section{Glossary}

\begin{itemize}[leftmargin=1.5em]
    \item \textbf{Coach Bot:} Automated system that sends welcome, summary, and nudge messages to keep users engaged.
    \item \textbf{Emoji Reactions:} Quick-feedback mechanism using emoji (e.g., \texttt{+1}, \texttt{heart}, \texttt{fire}).
    \item \textbf{Leaderboard:} Monthly ranking of channel members by check-in count.
    \item \textbf{Skeleton Loader:} Placeholder UI element shown while data is loading.
    \item \textbf{Toast:} Transient notification message that appears and dismisses automatically.
    \item \textbf{APScheduler:} Python library for scheduling background tasks at fixed intervals.
\end{itemize}

\end{document}
